% ===========================================================================
%  第 5 章 15 状态误差状态 EKF
% ===========================================================================
\section{15 状态误差状态 EKF}

本章建立主滤波器的完整状态空间模型：误差状态定义、连续时间系统矩阵
$\m{F}$ 的逐块推导、协方差传播、量测更新与误差注入。全部结果与
\file{ekf\_15\_state.h} 中 \code{PropagateStateFromIncrements}、
\code{InjectErrorState} 与 \code{MeasurementUpdate*} 系列逐行对照。

\subsection{误差状态定义}

固件以 \code{x\_}（15 维）存储误差状态，其索引布局为
（\file{ekf\_15\_state.h} 第 8--13 行、第 3891 行）：
\begin{equation}
  \dx = \begin{bmatrix}
    \delta\vp{p}^n \\ \delta\vp{v}^n \\ \delta\vp{\beta}^b \\
    \delta\vp{b}_a \\ \delta\vp{b}_g
  \end{bmatrix}
  = \begin{bmatrix}
    \text{NED 位置误差 (m)} \\
    \text{NED 速度误差 (m/s)} \\
    \text{体系姿态误差角 (rad)} \\
    \text{加速度计零偏误差 (m/s}^2\text{)} \\
    \text{陀螺零偏误差 (rad/s)}
  \end{bmatrix}.
  \label{eq:state-def}
\end{equation}
其中姿态误差采用\textbf{体系右乘}定义：
\begin{equation}
  q = q_{nom} \otimes \Exp(\delta\vp{\beta}^b),
  \label{eq:att-err-def}
\end{equation}
即真实姿态为标称姿态再右乘一个体系微旋转。\cref{eq:att-err-def} 与导航系左乘
定义 $q = \Exp(\delta\vp{\phi}^n)\otimes q_{nom}$ 不同（相差
$\m{C}_b^n$ 旋转），这一区别直接影响 $\m{F}$ 矩阵中姿态-速度耦合块的符号
（见第 5.3.4 节）。

\subsection{连续时间误差动力学}

对 \cref{eq:att-cont}--\cref{eq:pos-cont} 取变分，并保留与状态相关的所有一阶
项，得到误差动力学
\begin{align}
  \delta\dot{\vp{p}}^n &= \delta\vp{v}^n
    + \underbrace{\m{F}_{pp}\,\delta\vp{p}^n}_{\text{位置自耦合}},
  \label{eq:dp-dyn} \\
  \delta\dot{\vp{v}}^n &= -\sk{\m{C}_b^n\vp{f}^b}\,\delta\vp{\beta}^b
    - \m{C}_b^n\,\delta\vp{b}_a
    + \m{F}_{vv}\,\delta\vp{v}^n
    + \m{F}_{vp}\,\delta\vp{p}^n,
  \label{eq:dv-dyn} \\
  \delta\dot{\vp{\beta}}^b &= -\sk{\vp{\omega}_{ib}^b}\,\delta\vp{\beta}^b
    - \delta\vp{b}_g
    + \m{F}_{\beta v}\,\delta\vp{v}^n,
  \label{eq:dbeta-dyn} \\
  \delta\dot{\vp{b}}_a &= -\frac{1}{\tau_a}\delta\vp{b}_a + \vp{w}_a,
  \label{eq:dba-dyn} \\
  \delta\dot{\vp{b}}_g &= -\frac{1}{\tau_g}\delta\vp{b}_g + \vp{w}_g,
  \label{eq:dbg-dyn}
\end{align}
其中 $\vp{w}_a$、$\vp{w}_g$ 为高斯-马尔可夫驱动白噪声。各耦合矩阵在以下各节
推导。\cref{eq:dv-dyn} 中第一项来自比力项
$\delta(\m{C}_b^n\vp{f}^b) = \delta\m{C}_b^n\vp{f}^b$，而
$\delta\m{C}_b^n = \m{C}_b^n\sk{\delta\vp{\beta}^b}$（由 \cref{eq:small-rotation}
与 \cref{eq:skew-rot}），故
\begin{equation}
  \delta(\m{C}_b^n\vp{f}^b)
  = \m{C}_b^n\sk{\delta\vp{\beta}^b}\vp{f}^b
  = -\m{C}_b^n\sk{\vp{f}^b}\delta\vp{\beta}^b
  = -\sk{\m{C}_b^n\vp{f}^b}\,\delta\vp{\beta}^b.
  \label{eq:dv-att-coupling}
\end{equation}

\subsection{系统矩阵 F 的逐块推导}

将 \cref{eq:dp-dyn}--\cref{eq:dbg-dyn} 整理为
$\delta\dot{\vp{x}} = \m{F}\,\dx$，$\m{F}$ 的稀疏结构为
\begin{equation}
  \m{F} =
  \setcounter{MaxMatrixCols}{15}
  \begin{bmatrix}
    \m{F}_{pp} & \m{I}_3 & \m{0} & \m{0} & \m{0} \\
    \m{F}_{vp} & \m{F}_{vv} & -\sk{\m{C}_b^n\vp{f}^b} & -\m{C}_b^n & \m{0} \\
    \m{0} & \m{F}_{\beta v} & -\sk{\vp{\omega}_{ib}^b} & \m{0} & -\m{I}_3 \\
    \m{0} & \m{0} & \m{0} & -\tfrac{1}{\tau_a}\m{I}_3 & \m{0} \\
    \m{0} & \m{0} & \m{0} & \m{0} & -\tfrac{1}{\tau_g}\m{I}_3
  \end{bmatrix}.
  \label{eq:F-matrix}
\end{equation}

\subsubsection{位置-速度块 $\m{F}_{pv}$}

由 $\delta\dot{\vp{p}}^n = \delta\vp{v}^n$，$\m{F}(0{:}3,\ 3{:}6) = \m{I}_3$。
固件 \code{RefreshStateTransitionStaticBlocks} 第 2155 行静态写入。

\subsubsection{位置自耦合 $\m{F}_{pp}$}

位置误差 $\delta\vp{p}^n$ 的变化率对位置误差本身的偏导源于曲率半径与
$\cos\varphi$ 的非线性（对照 KF-GINS \cite{Farrell2000}）：
\begin{equation}
  \m{F}_{pp} =
  \begin{bmatrix}
    -\dfrac{v_D}{R_M+h} & 0 & \dfrac{v_N}{R_M+h} \\[1.2ex]
    \dfrac{v_E\tan\varphi}{R_N+h} & -\dfrac{v_D + v_N\tan\varphi}{R_N+h}
      & \dfrac{v_E}{R_N+h} \\[1.2ex]
    0 & 0 & 0
  \end{bmatrix}.
  \label{eq:Fpp}
\end{equation}
固件在完整 $\m{F}$ 矩阵路径（\code{BFS\_NAVIGATION\_EMBEDDED\_FULL\_F\_MATRIX=1}）
第 2371--2380 行写入 \cref{eq:Fpp} 的五个非零元素。

\subsubsection{速度-位置块 $\m{F}_{vp}$ 与重力梯度}

NED 速度误差动力学中包含位置误差引起的重力与曲率变化：
\begin{equation}
  \m{F}_{vp} = \begin{bmatrix}
    0 & 0 & 0 \\ 0 & 0 & 0 \\ 0 & 0 & \dfrac{2g}{R_{mean}}
  \end{bmatrix},
  \label{eq:Fvp}
\end{equation}
其中 $F(5,2) = 2g/R_{mean}$ 为\textbf{高度通道重力梯度}。推导：由
\cref{eq:gravity-alt}，$\partial g/\partial h = -2g/R_{mean}$（对 $h\ll R$）；
NED 地向误差 $\delta p_D = -\delta h$ 增大（下降）时，重力矢量的地向分量增大，
故
\begin{equation}
  \delta\dot{v}_D \supset \frac{\partial g}{\partial h}\,\delta h
  = \left(-\frac{2g}{R_{mean}}\right)(-\delta p_D)
  = \frac{2g}{R_{mean}}\,\delta p_D.
  \label{eq:f52}
\end{equation}
固件第 2322 行写入 $F(5,2) = 2.0\cdot g / mean\_radius$。

\subsubsection{速度-姿态块 $\m{F}_{v\beta}$}

由 \cref{eq:dv-att-coupling}，
\begin{equation}
  \m{F}_{v\beta} = -\sk{\m{C}_b^n\vp{f}^b},
  \label{eq:Fvbeta}
\end{equation}
固件第 2323 行以 \code{-t\_b2ned\_mid * Skew(ins\_accel\_mps2\_)} 实现。

\subsubsection{速度-加计零偏块 $\m{F}_{v b_a}$}

\begin{equation}
  \m{F}_{v b_a} = -\m{C}_b^n,
  \label{eq:Fvba}
\end{equation}
固件第 2324 行写入。

\subsubsection{速度自耦合 $\m{F}_{vv}$（哥氏/传输速率）}

速度误差之间的耦合来自科氏项 $-(2\vp{\omega}_{ie}^n +
\vp{\omega}_{en}^n)\cross\vp{v}^n$ 对速度的偏导（对照 KF-GINS）：
\begin{equation}
  \m{F}_{vv} =
  \begin{bmatrix}
    \dfrac{v_D}{R_M+h} &
    -2\left(\omega_e\sin\varphi + \dfrac{v_E\tan\varphi}{R_N+h}\right) &
    \dfrac{v_N}{R_M+h} \\[1.2ex]
    2\omega_e\sin\varphi + \dfrac{v_E\tan\varphi}{R_N+h} &
    \dfrac{v_D + v_N\tan\varphi}{R_N+h} &
    2\omega_e\cos\varphi + \dfrac{v_E}{R_N+h} \\[1.2ex]
    -\dfrac{2v_N}{R_M+h} &
    -2\left(\omega_e\cos\varphi + \dfrac{v_E}{R_N+h}\right) & 0
  \end{bmatrix}.
  \label{eq:Fvv}
\end{equation}
固件第 2352--2366 行逐元素写入 \cref{eq:Fvv}（九个元素），注释明确其为
``哥氏力/传输速率自耦合''。

\subsubsection{姿态自旋块 $\m{F}_{\beta\beta}$}

\begin{equation}
  \m{F}_{\beta\beta} = -\sk{\vp{\omega}_{ib}^b},
  \label{eq:Fbb}
\end{equation}
固件第 2325 行写入 \code{-Skew(ins\_gyro\_radps\_)}。

\subsubsection{姿态-陀螺零偏块 $\m{F}_{\beta b_g}$}

由 $\delta\dot{\vp{\beta}}^b = -\delta\vp{b}_g$，
\begin{equation}
  \m{F}_{\beta b_g} = -\m{I}_3,
  \label{eq:Fbbg}
\end{equation}
固件第 2156 行静态写入。

\subsubsection{姿态-速度块 $\m{F}_{\beta v}$（传输速率雅可比）}

这是\textbf{误差状态定义相关}的关键块。体系右乘误差 $\delta\vp{\beta}^b$
的动力学中包含传输速率项：导航系转动 $\vp{\omega}_{en}^n$ 的变化
$\delta\vp{\omega}_{en}^n$ 经 $\m{C}_n^b$ 旋转到机体后以负号进入
$\delta\dot{\vp{\beta}}^b$：
\begin{equation}
  \delta\dot{\vp{\beta}}^b \supset -\m{C}_n^b\,\delta\vp{\omega}_{en}^n
  = -\m{C}_n^b\,\m{M}\,\delta\vp{v}^n,
  \label{eq:dbeta-transport}
\end{equation}
其中 $\m{M}$ 为传输速率雅可比（由 \cref{eq:nav-rate} 对速度求导）：
\begin{equation}
  \m{M} =
  \begin{bmatrix}
    0 & \dfrac{1}{R_N+h} & 0 \\[1.2ex]
    -\dfrac{1}{R_M+h} & 0 & 0 \\[1.2ex]
    0 & -\dfrac{\tan\varphi}{R_N+h} & 0
  \end{bmatrix}.
  \label{eq:transport-jacobian}
\end{equation}
故
\begin{equation}
  \m{F}_{\beta v} = -\m{C}_n^b\,\m{M} = -\m{C}_n^b\,\m{M}.
  \label{eq:Fbv}
\end{equation}
固件第 2397--2401 行按 \cref{eq:Fbv} 实现：
\code{fs\_.block<3,3>(6,3) = -t\_b2ned\_mid.transpose() * transport\_jacobian}。
\cref{tab:Fbv-compare} 给出与导航系左乘误差定义的对照。

\begin{table}[H]
\centering
\caption{姿态误差定义对 $\m{F}_{\beta v}$ 的影响}
\label{tab:Fbv-compare}
\small
\begin{tabularx}{\textwidth}{@{}L{4.2cm}L{4.8cm}X@{}}
\toprule
\textbf{误差定义} & $\m{F}_{\beta v}$ & \textbf{说明} \\
\midrule
体系右乘（本固件） & $-\m{C}_n^b\m{M}$ & 含旋转项，水平时 $\m{C}_n^b=\m{I}$ 退化为 $-\m{M}$ \\
导航系左乘（KF-GINS） & $+\m{M}$ & 无旋转项 \\
\bottomrule
\end{tabularx}
\end{table}

\begin{remark}
固件注释（第 2393--2395 行）明确指出：当姿态水平（$\m{C}_b^n = \m{I}$）时，
两者差一个负号。该差异源于误差角的坐标化（body vs.\ nav），两种定义在数学上
严格同构——主机回归测试 \file{test/ekf\_host\_regression.cpp} 专门验证了
$F(\Phi,V) = -\m{C}_n^b\m{M}$ 的符号与坐标一致性。
\end{remark}

\subsubsection{零偏衰减块}

\begin{equation}
  \m{F}_{b_a b_a} = -\frac{1}{\tau_a}\m{I}_3, \qquad
  \m{F}_{b_g b_g} = -\frac{1}{\tau_g}\m{I}_3,
  \label{eq:F-bias-decay}
\end{equation}
固件第 2157--2158 行写入 \code{accel\_markov\_bias\_} 与
\code{gyro\_markov\_bias\_}（初始化于第 575--578 行）。

\subsection{简化 F 矩阵与完整 F 矩阵}

\cref{eq:F-matrix} 所示结构在固件中通过编译宏
\code{BFS\_NAVIGATION\_EMBEDDED\_FULL\_F\_MATRIX} 控制：
\begin{itemize}
  \item \textbf{简化版}（宏=0）：仅包含核心块 $\m{F}_{pv}$、
    $\m{F}_{v\beta}$、$\m{F}_{v b_a}$、$\m{F}_{\beta\beta}$、
    $\m{F}_{\beta b_g}$、$F(5,2)$ 与零偏衰减块；
  \item \textbf{完整版}（宏=1，固件默认）：额外包含 $\m{F}_{vv}$（9 项）、
    $\m{F}_{pp}$（5 项）、$\m{F}_{\beta v}$（3 项）。
\end{itemize}
固件第 103--118 行的注释给出适用边界：短时飞行（$<$2\,min）+ MEMS IMU + 高频
GNSS 修正时简化版已足够（升级收益 $<$0.1°）；长航时（$>$5\,min）或高速
（$>$500\,m/s）建议开启完整版。200\,Hz 嵌入式路径下完整版额外开销约 15 次
float 乘加，可忽略。

\subsection{可观测性分析}

误差状态的\textbf{可观测性}决定哪些状态能被量测组合收敛，直接解释第 7 章静止
辅助设计与第 10 章仿真 S1 的结果。对线性时变系统
\begin{equation}
  \delta\dot{\vp{x}} = \m{F}(t)\dx, \qquad
  \vp{y} = \m{H}(t)\dx,
  \label{eq:obs-sys}
\end{equation}
其可观测性可由\textbf{可观测性格拉姆矩阵}刻画：
\begin{equation}
  \m{W}_O(t_0, t_1) = \int_{t_0}^{t_1}
    \bm{\Phi}\T(s, t_0)\,\m{H}\T(s)\,\m{H}(s)\,\bm{\Phi}(s, t_0)\,ds,
  \label{eq:gramian}
\end{equation}
$\m{W}_O$ 满秩当且仅当系统在 $[t_0, t_1]$ 上完全可观 \cite{crassidis2011}。

\subsubsection{静止工况（无 GNSS）}

静止时 $\m{C}_b^n = \m{I}$、$\vp{f}^b = (0,0,g)$、$\vp{\omega} = \m{0}$。
忽略地球项（量级 $\le 10^{-5}$），$\m{F}$ 化简为
\begin{equation}
  \m{F} = \begin{bmatrix}
    \m{0} & \m{I}_3 & \m{0} & \m{0} & \m{0} \\
    \m{0} & \m{0} & -\sk{\vp{g}^b} & -\m{I}_3 & \m{0} \\
    \m{0} & \m{0} & \m{0} & \m{0} & -\m{I}_3 \\
    \m{0} & \m{0} & \m{0} & \m{0} & \m{0} \\
    \m{0} & \m{0} & \m{0} & \m{0} & \m{0}
  \end{bmatrix},
  \label{eq:F-static}
\end{equation}
其中 $\vp{g}^b = (0,0,g)$。静止辅助量测组合为
$\vp{y} = \left[\vp{v}^n,\ \m{C}_n^b\vp{g}_{sf}^n,\ \vp{\omega}_{ib}^b\right]$。
代入 \cref{eq:gramian} 逐块检验可观测子空间，结论如下：
\begin{itemize}
  \item \textbf{姿态-速度-水平零偏的耦合}：速度误差动力学
    $\delta\dot{\vp{v}} = -\sk{\vp{g}^b}\delta\vp{\beta}^b - \delta\vp{b}_a$，
    静止时 $\sk{\vp{g}^b} = g\begin{bmatrix} 0&1&0 \\ -1&0&0 \\ 0&0&0 \end{bmatrix}$
    的零空间为 $\operatorname{span}\{(0,0,1)\}$——即 yaw 误差 $\delta\beta_D$
    与 $\delta\vp{v}$ 解耦；而 $\delta\beta_{N/E}$ 与 $\delta b_{a,N/E}$
    通过同一方程耦合，两者之和（在 $\delta\vp{v}$ 观测下）可观、之差不可观。
    这正是第 10.2 节仿真中水平加速度零偏不收敛的数学根源；
  \item \textbf{垂直零偏可观}：$\delta b_{a,D}$ 与 $\delta\dot{v}_D$ 直接
    耦合且无倾角竞争项（$\sk{\vp{g}^b}$ 的第三行全零），ZUPT 的
    $v_D = 0$ 观测使其可观；
  \item \textbf{陀螺零偏三轴可观}：StaticGyro 直接观测
    $\delta\vp{b}_g$（\cref{eq:H-R-sgyro} 的 $\m{H}_{sgyro} = [\m{0}\ \m{I}_3]$），
    三轴均直接可观；
  \item \textbf{yaw 不可观}：$\delta\beta_D$ 不进入任何静止量测方程
    （重力方向 $\m{g}_{sf}^n = (0,0,-1)$ 经 $\m{C}_n^b$ 旋转到机体后仅含
    roll/pitch 信息），且与速度/零偏解耦，故完全不可观——初始协方差
    $\pi$ 保持。
\end{itemize}

\subsubsection{运动工况（GNSS 有效）}

GNSS 六维量测 $\m{H}_{gps} = [\m{I}_6\ \m{0}]$ 直接观测位置与速度。由
\cref{eq:F-static} 的耦合结构，速度观测经 $\m{F}_{v\beta}$、
$\m{F}_{v b_a}$ 与 $\m{F}_{\beta b_g}$ 的积分链路间接观测姿态与零偏。若
$\m{C}_b^n\vp{f}^b$ 保持非零（有机动比力），$\sk{\m{C}_b^n\vp{f}^b}$ 提供
姿态-速度耦合，使 $\delta\vp{\beta}^b$ 可观；若飞行全程比力方向不变（匀速
直线），耦合恒定，姿态的可观测性退化为零偏与倾角的组合——这与静止情形同构。
\textbf{激励越丰富（转弯/加减速），姿态与零偏的可观性越强}，这是组合导航
``需要机动激励来标定传感器''的经验定律的数学表达
\cite{Groves2013}。

\begin{remark}
可观测性分析解释了固件的两个关键设计：一是静止辅助必须三量测配合
（ZUPT 约束速度、Gravity 约束倾角、StaticGyro 约束陀螺零偏）才能收敛可观的
子空间；二是水平加速度零偏不可观是\textbf{系统固有属性}，只能通过 GNSS
运动激励解耦——固件在 GNSS 有效时的动态权重与高频位置/速度更新正是为了
持续激励该耦合链。
\end{remark}

\subsection{协方差传播}

\subsubsection{状态转移矩阵}

由第 2.6 节的二阶截断，离散状态转移矩阵为
\cref{eq:phi-second}。固件第 2424--2425 行实现
\code{phi\_ = I + fs\_dt + 0.5*fs\_dt*fs\_dt}。

\subsubsection{连续过程噪声矩阵}

噪声输入矩阵 $\m{G}_s$（15$\times$12，第 2443--2446 行）为
\begin{equation}
  \m{G}_s = \begin{bmatrix}
    \m{0} & \m{0} & \m{0} & \m{0} \\
    -\m{C}_b^n & \m{0} & \m{0} & \m{0} \\
    \m{0} & -\m{I}_3 & \m{0} & \m{0} \\
    \m{0} & \m{0} & \m{I}_3 & \m{0} \\
    \m{0} & \m{0} & \m{0} & \m{I}_3
  \end{bmatrix},
  \label{eq:Gs}
\end{equation}
连续噪声谱密度 $\m{R}_w$（12$\times$12，第 513--530 行）为
\begin{equation}
  \m{R}_w = \diag\left(\sigma_{aw}^2\m{I}_3,\ \sigma_{gw}^2\m{I}_3,\
    \frac{2\sigma_{aB}^2}{\tau_a}\m{I}_3,\ \frac{2\sigma_{gB}^2}{\tau_g}\m{I}_3\right),
  \label{eq:Rw}
\end{equation}
其中各项含义与固件成员变量对照见 \cref{tab:noise-params}。

\begin{table}[H]
\centering
\caption{过程噪声参数与固件成员对照（\file{ekf\_15\_state.h} 第 3743--3785 行）}
\label{tab:noise-params}
\small
\begin{tabularx}{\textwidth}{@{}L{3.4cm}C{3.2cm}C{3.2cm}X@{}}
\toprule
\textbf{物理量} & \textbf{符号} & \textbf{固件默认/运行值} & \textbf{固件成员} \\
\midrule
加速度计白噪声 & $\sigma_{aw}$ & 0.05 / 0.25 m/s² & \code{accel\_std\_mps2} \\
陀螺白噪声 & $\sigma_{gw}$ & 0.00175 / 0.0015 rad/s & \code{gyro\_std\_radps} \\
加计零偏 Markov 驱动 & $\sigma_{aB}$ & 0.01 m/s² & \code{accel\_markov\_bias\_std} \\
陀螺零偏 Markov 驱动 & $\sigma_{gB}$ & 0.00025 / 0.00524 rad/s & \code{gyro\_markov\_bias\_std} \\
加计零偏相关时间 & $\tau_a$ & 100 s & \code{accel\_tau\_s} \\
陀螺零偏相关时间 & $\tau_g$ & 50 / 180 s & \code{gyro\_tau\_s} \\
\bottomrule
\end{tabularx}
\end{table}

\begin{remark}
固件运行配置（\file{navigation\_task.cpp} 第 846--867 行）针对 ICM-42688 与
共轴双桨振动环境做了专门调整：陀螺白噪声收紧至 0.0015\,rad/s（信任陀螺），
陀螺零偏 Markov 驱动放宽至 0.00524\,rad/s（0.3\,dps），加计白噪声放宽至
0.25\,m/s²（抵抗震动）。注释记录：0.05\,dps 的过小驱动噪声会导致 P 矩阵坍缩、
卡尔曼增益趋零、Z 轴零偏无法修正。
\end{remark}

\subsubsection{离散过程噪声与协方差递推}

默认高精度路径（第 2552--2603 行）使用双精度计算：
\begin{align}
  \m{Q}_d &= \frac{dt}{6}\left(\m{Q}_c + 4\bm{\Phi}_{\frac12}\m{Q}_c
    \bm{\Phi}_{\frac12}^\top + \bm{\Phi}\m{Q}_c\bm{\Phi}^\top\right),
  \label{eq:qd-path} \\
  \m{P}_{k} &= \bm{\Phi}\m{P}_{k-1}\bm{\Phi}^\top + \m{Q}_d,
  \label{eq:p-prop}
\end{align}
其中 $\m{Q}_c = \m{G}_s\m{R}_w\m{G}_s\T$。固件第 2601--2603 行以
\code{double} 中间量执行 $\bm{\Phi}\m{P}\bm{\Phi}^\top$ 再转回 \code{float}，
降低 15$\times$15 矩阵乘法在高动态大协方差下的舍入误差。

快速路径（\code{BFS\_NAVIGATION\_EMBEDDED\_FAST\_PROPAGATION}，默认关闭）使用
float 一阶近似 $\m{P} \approx \m{P} + \m{F}dt\m{P} + \m{P}\m{F}^\top dt + \m{Q}_d$，
并利用 $\m{Q}_d$ 的块对角结构直写对角线（第 2481--2541 行）。固件注释明确其
仅用于 ESP32 桌面 bring-up，不适用于飞控闭环。

\subsection{量测更新与 Joseph 协方差}

各量测更新（GNSS、速度、重力、姿态、航向、气压高度）统一遵循第 2.7 节的
ES-EKF 递推 \cref{eq:es-innovation}--\cref{eq:es-update}，协方差均采用
Joseph 形式。以 GNSS 六维量测为例（\file{ekf\_15\_state.h} 第 959--972 行）：
\begin{align}
  \m{K} &= \m{P}\m{H}^\top\m{S}\inv, &
  \m{S} &= \m{H}\m{P}\m{H}^\top + \m{R}, \label{eq:kgain} \\
  \m{P} &= \left(\m{I}-\m{K}\m{H}\right)\m{P}
    \left(\m{I}-\m{K}\m{H}\right)\T + \m{K}\m{R}\m{K}^\top.
  \label{eq:joseph}
\end{align}
其中 $\m{H}$、$\m{R}$ 的具体形式见第 6 章。固件对 $S$ 矩阵求逆使用
LDLT 分解（\code{InvertSymmetricPositiveDefiniteMatrix}，第 1550--1567 行）
并在反演前后做有限值检查，病态时拒绝该帧量测。

\subsection{误差状态注入}

量测更新得到 $\dx = \m{K}\vp{y}$ 后，\code{InjectErrorState}
（第 2623--2661 行）将误差注入全量状态：
\begin{enumerate}
  \item \textbf{位置}：按 WGS-84 几何把 NED 误差换算回经纬高
    （第 2627--2639 行）：
    \begin{align}
      \delta\varphi &= \frac{\dx(0)}{R_M + h}, &
      \delta\lambda &= \frac{\dx(1)}{(R_N + h)\cos\varphi}, &
      \delta h &= -\dx(2);
      \label{eq:inject-pos}
    \end{align}
  \item \textbf{速度}：$\vp{v}^n \leftarrow \vp{v}^n + \dx(3{:}6)$;
  \item \textbf{姿态}：采用指数映射而非一阶近似（第 2643--2651 行）：
    \begin{equation}
      q \leftarrow \left(q\otimes\Exp\left(\dx(6{:}9)\right)\right)_{norm},
      \label{eq:inject-att}
    \end{equation}
    并在 $q_w<0$ 时整体取反；
  \item \textbf{零偏}：$\vp{b}_a \leftarrow \vp{b}_a + \dx(9{:}12)$，
    $\vp{b}_g \leftarrow \vp{b}_g + \dx(12{:}15)$，同时回退本帧去偏比力/角速度
    （第 2654--2658 行）；
  \item \textbf{清零}：$\dx \leftarrow \m{0}$（第 2660 行）。
\end{enumerate}

\begin{remark}
\cref{eq:inject-pos} 的经纬度换算使用更新前的纬度与高度计算比例尺，避免
分量顺序引入二阶耦合误差（第 2634 行注释）。位置误差以 \code{double} 累积在
LLA 中，只有相对 NED 输出时才转换回 \code{float}，保证绝对位置精度。
\end{remark}

\subsection{协方差稳定化}

\code{StabilizeCovariance}（第 1803--1899 行）采用三级递进的数值稳定策略：
\begin{enumerate}
  \item \textbf{对称化} + Gershgorin 下界检查
    （\code{SymmetrizeAndComputeGershgorinLowerBound}，第 1984--2023 行）：
    若下界 $\ge -\varepsilon_{tol}$，直接判定 PSD 并返回（多数健康帧走此路径）；
  \item \textbf{LDLT 正定性检查}：Gershgorin 下界为负并不代表矩阵损坏（判据
    保守），LDLT 通过则返回（第 1858--1871 行）；
  \item \textbf{特征值投影}：对负特征值 $\lambda_i < \max(\lambda_{min},
    \lambda_{max}\cdot 10^{-10})$ 投影到该下界后重构
    $\m{P} = \m{V}\diag(\lambda_i)\m{V}\T$（第 1872--1898 行）。
\end{enumerate}
三级策略将昂贵的 15$\times$15 特征分解从热路径移出，仅在异常帧触发。运行期
稳定化分频（\code{propagation\_stabilize\_divider}=8）使完整稳定化以
25\,Hz 执行（\file{navigation\_task.cpp} 第 872--873 行）。
